% ═══════════════════════════════════════════════════════════════════════════
% UNTERRICHTSLANGENTWURF: Me gusta - Vorlieben ausdrücken
% Fach: Spanisch (Oberstufe Spätbeginner)
% Gemäß: Fächerübergreifende Handreichung M-V (ZLB/IQ M-V)
% ═══════════════════════════════════════════════════════════════════════════

\documentclass[11pt, a4paper]{article}

% ═══════════════════════════════════════════════════════════════════════════
% PAKETE
% ═══════════════════════════════════════════════════════════════════════════

\usepackage[utf8]{inputenc}
\usepackage[T1]{fontenc}
\usepackage[ngerman]{babel}
\usepackage{helvet}
\renewcommand{\familydefault}{\sfdefault}

\usepackage[margin=2.5cm, left=3.5cm, top=2.5cm, bottom=2.5cm]{geometry}
\usepackage{setspace}
\onehalfspacing

\usepackage{enumitem}
\usepackage{tabularx}
\usepackage{booktabs}
\usepackage{array}
\usepackage{longtable}
\usepackage{multirow}

\usepackage{xcolor}
\usepackage{tcolorbox}
\usepackage{fancyhdr}
\usepackage{lastpage}
\usepackage{titlesec}
\usepackage{parskip}

% Hyperlinks (für Inhaltsverzeichnis)
\usepackage[hidelinks]{hyperref}

% ═══════════════════════════════════════════════════════════════════════════
% FARBEN
% ═══════════════════════════════════════════════════════════════════════════

\definecolor{spanisch}{HTML}{c41e3a}
\definecolor{grau}{HTML}{666666}
\definecolor{hellgrau}{HTML}{f5f5f5}
\definecolor{success}{HTML}{2a7a4b}

% ═══════════════════════════════════════════════════════════════════════════
% ÜBERSCHRIFTEN
% ═══════════════════════════════════════════════════════════════════════════

\titleformat{\section}
    {\normalfont\large\bfseries\color{spanisch}}
    {\thesection}{1em}{}

\titleformat{\subsection}
    {\normalfont\normalsize\bfseries}
    {\thesubsection}{1em}{}

\titlespacing*{\section}{0pt}{18pt}{10pt}
\titlespacing*{\subsection}{0pt}{12pt}{6pt}

% ═══════════════════════════════════════════════════════════════════════════
% KOPF- UND FUSSZEILE
% ═══════════════════════════════════════════════════════════════════════════

\pagestyle{fancy}
\fancyhf{}
\renewcommand{\headrulewidth}{0.4pt}
\renewcommand{\footrulewidth}{0pt}

\lhead{\small Unterrichtslangentwurf -- Spanisch}
\rhead{\small Me gusta -- Vorlieben ausdrücken}
\cfoot{\small Seite \thepage\ von \pageref{LastPage}}

% ═══════════════════════════════════════════════════════════════════════════
% BEFEHLE
% ═══════════════════════════════════════════════════════════════════════════

% Spanisch-Highlight
\newcommand{\es}[1]{\textcolor{spanisch}{\textit{#1}}}

% Info-Box
\newtcolorbox{infobox}[1][]{
    colback=hellgrau,
    colframe=spanisch,
    boxrule=0.8pt,
    arc=0pt,
    left=8pt, right=8pt, top=6pt, bottom=6pt,
    #1
}

% ═══════════════════════════════════════════════════════════════════════════
% DOKUMENT
% ═══════════════════════════════════════════════════════════════════════════

\begin{document}

% ═══════════════════════════════════════════════════════════════════════════
% DECKBLATT (Bestandteil 1 + 2)
% ═══════════════════════════════════════════════════════════════════════════

\thispagestyle{empty}

\begin{center}
    \vspace*{2cm}

    {\Large\bfseries Unterrichtslangentwurf}

    \vspace{0.5cm}

    {\large Fach: Spanisch}

    \vspace{2cm}

    {\LARGE\bfseries\color{spanisch} Me gusta}

    \vspace{0.3cm}

    {\large Vorlieben ausdrücken mit dem Verb \es{gustar}}

    \vspace{3cm}

    \begin{tabular}{ll}
        \textbf{Referendar/in:} & [Name] \\[6pt]
        \textbf{Schule:} & \underline{\hspace{5cm}} \\[6pt]
        \textbf{Klasse:} & Einführungsphase (Jg. 11), Spanisch Spätbeginner \\[6pt]
        \textbf{Datum:} & [Datum] \\[6pt]
        \textbf{Zeit:} & [Uhrzeit], 45 Minuten \\[6pt]
        \textbf{Raum:} & [Raum] \\[6pt]
    \end{tabular}

    \vspace{2cm}

    \begin{tabular}{ll}
        \textbf{Mentor/in:} & [Name] \\[6pt]
        \textbf{Fachleiter/in:} & [Name] \\[6pt]
        \textbf{Studienleiter/in:} & [Name] \\[6pt]
        \textbf{Schulleiter/in:} & [Name] \\[6pt]
    \end{tabular}

    \vfill

    \hrule
    \vspace{0.3cm}
    {\small\color{grau} Unterrichtseinheit: Gustos -- Hobbys und Freizeit (Stunde 1 von 3)}

\end{center}

\newpage

% ═══════════════════════════════════════════════════════════════════════════
% INHALTSVERZEICHNIS
% ═══════════════════════════════════════════════════════════════════════════

\tableofcontents
\newpage

% ═══════════════════════════════════════════════════════════════════════════
% 3. BEDINGUNGSANALYSE
% ═══════════════════════════════════════════════════════════════════════════

\section{Bedingungsanalyse}

\subsection{Statistik der Lerngruppe}

Die Lerngruppe besteht aus 18 Schülerinnen und Schülern (11 weiblich, 7 männlich) der Einführungsphase (Jahrgangsstufe 11). Es handelt sich um einen Kurs für Spätbeginner, d.\,h. alle Lernenden haben Spanisch erst in Klasse 11 als neu einsetzende Fremdsprache gewählt. Der Kurs findet dreistündig statt.

\subsection{Psychologische Besonderheiten und Interessensituation}

Die Lerngruppe zeigt eine überwiegend positive Einstellung zum Fach Spanisch. Die Mehrheit der Schülerinnen und Schüler hat das Fach aus Interesse an der spanischsprachigen Welt (Reisen, Musik, Kultur) gewählt. Das Lernklima ist kooperativ; Partnerarbeit funktioniert gut.

Die Motivation ist zu Beginn des Schuljahres hoch, da viele SuS erstmals eine neue Sprache „von Null" lernen und schnelle Fortschritte erleben. Einige SuS zeigen Unsicherheit bei grammatischen Strukturen, die vom Deutschen oder Englischen abweichen.

\subsection{Individuelle Besonderheiten}

\begin{itemize}[leftmargin=2em]
    \item Ein Schüler (S1) hat diagnostizierte LRS; er erhält Arbeitsblätter in größerer Schrift und mehr Zeit bei schriftlichen Aufgaben.
    \item Eine Schülerin (S2) hat Spanisch-Vorkenntnisse durch familiären Hintergrund (Mutter aus Chile); sie wird als Expertin bei Partnerarbeit eingesetzt.
    \item Zwei SuS sind zurückhaltend bei mündlicher Beteiligung; sie werden gezielt in Partnerarbeit aktiviert.
\end{itemize}

\subsection{Arbeitsvoraussetzungen}

\textbf{Materielle Voraussetzungen:}
\begin{itemize}[leftmargin=2em]
    \item Der Unterrichtsraum verfügt über Beamer, interaktives Whiteboard und WLAN.
    \item Die SuS können Tablets (1:1-Ausstattung) oder eigene Smartphones nutzen.
    \item Arbeitsblätter werden analog ausgeteilt (Hefter-Prinzip).
\end{itemize}

\textbf{Methodische Voraussetzungen:}
\begin{itemize}[leftmargin=2em]
    \item Die SuS sind mit digitalen Lernpfaden vertraut (Einführung in Woche 1).
    \item Partnerarbeit und Think-Pair-Share sind etablierte Arbeitsformen.
    \item Grammatische Fachbegriffe (Subjekt, Verb, Pronomen) sind aus dem Englisch- und Deutschunterricht bekannt.
\end{itemize}

\subsection{Fachbezogener Entwicklungsstand}

Die Lerngruppe befindet sich im ersten Lernjahr (Einführungsphase) auf dem Weg zu GER-Niveau A2. Bisherige Unterrichtsinhalte:

\begin{itemize}[leftmargin=2em]
    \item Sich vorstellen, Herkunft und Wohnort angeben
    \item Verben auf -ar, -er, -ir im Präsens (regelmäßig)
    \item Artikel und Genus, Plural
    \item Possessivpronomen (mi, tu, su...)
    \item Zahlen 1-100, Datum, Uhrzeit
\end{itemize}

Die \es{gustar}-Konstruktion ist \textbf{neu}. Die SuS kennen jedoch indirekte Objektpronomen noch nicht systematisch. Die Stunde führt diese im Kontext von \es{gustar} ein.

\subsection{Erfahrungen mit der Lerngruppe}

Die Lerngruppe ist mir seit Schuljahresbeginn (12 Wochen) bekannt. Die Leistungsfähigkeit ist heterogen: etwa ein Drittel zeigt überdurchschnittliche Motivation und schnelle Auffassung, ein weiteres Drittel arbeitet solide mit, und das letzte Drittel benötigt mehr Unterstützung und Zeit.

% ═══════════════════════════════════════════════════════════════════════════
% 4. SACHANALYSE
% ═══════════════════════════════════════════════════════════════════════════

\section{Sachanalyse}

\subsection{Das Verb \es{gustar} -- linguistische Einordnung}

Das spanische Verb \es{gustar} gehört zur Gruppe der sogenannten \textbf{„verba sentiendi"} oder \textbf{psychologischen Verben} (verbos de afección), bei denen das Experiencer-Argument (die Person, die etwas empfindet) nicht als grammatisches Subjekt, sondern als indirektes Objekt realisiert wird (vgl. Bosque \& Demonte, 1999, Bd. 2, Kap. 24).

\textbf{Syntaktische Struktur:}
\begin{center}
    \textit{(A + Dativobjekt)} + \textbf{Dativpronomen} + \es{gustar} + \textbf{Subjekt}
\end{center}

Beispiel: \es{A María le gustan los libros.}

\begin{itemize}[leftmargin=2em]
    \item \textit{A María}: präpositionale Phrase zur Verdeutlichung des Experiencers
    \item \textit{le}: indirektes Objektpronomen (obligatorisch)
    \item \textit{gustan}: Verb, kongruiert mit dem grammatischen Subjekt
    \item \textit{los libros}: grammatisches Subjekt (bestimmt Numerus des Verbs)
\end{itemize}

\subsection{Kontrastive Betrachtung: Spanisch -- Deutsch}

Im Deutschen existiert eine analoge Struktur mit dem Verb „gefallen":

\begin{center}
\begin{tabular}{ll}
    \textbf{Deutsch:} & \underline{Mir} gefällt \underline{die Musik}. \\
    \textbf{Spanisch:} & \underline{Me} gusta \underline{la música}. \\
\end{tabular}
\end{center}

Diese Parallele ist didaktisch wertvoll: Deutsche Muttersprachler können die \es{gustar}-Struktur über die Brücke „gefallen" erschließen. Der häufige Interferenzfehler („*Yo gusto...") entsteht durch Transfer vom englischen „I like...", bei dem der Experiencer Subjekt ist.

\subsection{Semantik und Verwendung}

\es{Gustar} drückt ein positives affektives Verhältnis aus. Es kann mit Nomen (Singular/Plural) oder Infinitiven kombiniert werden:

\begin{itemize}[leftmargin=2em]
    \item Mit Nomen: \es{Me gusta el café.} / \es{Me gustan los gatos.}
    \item Mit Infinitiv: \es{Me gusta bailar.} (immer Singular \es{gusta})
    \item Mit mehreren Infinitiven: \es{Me gusta leer y escribir.} (Singular, wenn als Einheit)
\end{itemize}

\subsection{Die indirekten Objektpronomen}

\begin{center}
\begin{tabular}{|c|c|c|}
    \hline
    \textbf{Person} & \textbf{Pronomen} & \textbf{Optionale Verstärkung} \\
    \hline
    1. Sg. & me & a mí \\
    2. Sg. & te & a ti \\
    3. Sg. & le & a él / a ella / a usted \\
    1. Pl. & nos & a nosotros/-as \\
    2. Pl. & os & a vosotros/-as \\
    3. Pl. & les & a ellos / a ellas / a ustedes \\
    \hline
\end{tabular}
\end{center}

Die präpositionalen Formen (\es{a mí, a ti...}) dienen der \textbf{Emphase} oder \textbf{Disambiguierung} bei \es{le/les}.

\subsection{Weitere Verben mit gleicher Struktur}

Neben \es{gustar} folgen weitere Verben diesem Muster:
\begin{itemize}[leftmargin=2em]
    \item \es{encantar} (begeistern): \es{Me encanta viajar.}
    \item \es{interesar} (interessieren): \es{¿Te interesa el arte?}
    \item \es{molestar} (stören): \es{Me molesta el ruido.}
    \item \es{doler} (schmerzen): \es{Me duele la cabeza.}
\end{itemize}

Diese werden in Folgestunden thematisiert; die Struktur wird hier grundlegend eingeführt.

% ═══════════════════════════════════════════════════════════════════════════
% 5. DIDAKTISCHE ÜBERLEGUNGEN
% ═══════════════════════════════════════════════════════════════════════════

\section{Didaktische Überlegungen}

\subsection{Stellung der Stunde in der Unterrichtseinheit}

Die vorliegende Stunde ist die \textbf{erste von drei Stunden} der Unterrichtseinheit „Gustos -- Hobbys und Freizeit". Die Sequenz ist wie folgt geplant:

\begin{enumerate}[leftmargin=2em]
    \item \textbf{Diese Stunde:} Einführung \es{gustar} + Grundwortschatz Hobbys
    \item Stunde 2: Vertiefung (encantar, interesar) + dialogisches Sprechen
    \item Stunde 3: Anwendung: Gespräch über Vorlieben + schriftliche Produktion
\end{enumerate}

\textbf{Funktion der Stunde:} Einführungsstunde mit dem Ziel, die ungewohnte Satzstruktur zu verstehen und erste Anwendungssicherheit zu erlangen.

\subsection{Ziele der Unterrichtseinheit}

\begin{enumerate}[leftmargin=2em]
    \item Die SuS verstehen die Struktur psychologischer Verben wie \es{gustar}.
    \item Die SuS können Vorlieben und Abneigungen ausdrücken und erfragen.
    \item Die SuS erweitern ihren Wortschatz im Themenfeld Freizeit/Hobbys.
    \item Die SuS können kurze Gespräche über Vorlieben führen (GER A2).
\end{enumerate}

\subsection{Legitimation}

\textbf{Bezug zum Rahmenplan:}

Der Rahmenplan Spanisch für die gymnasiale Oberstufe M-V (2019) fordert im Bereich „Verfügen über sprachliche Mittel" (Grammatik):
\begin{quote}
    „Die Schülerinnen und Schüler verwenden ein gefestigtes Repertoire häufig verwendeter grammatischer Strukturen [...] und können elementare Mittel zur Realisierung von Sprechabsichten einsetzen." (RP Spa Sek II, S. 12)
\end{quote}

Das Verb \es{gustar} ist ein \textbf{hochfrequentes Alltagsverb}, das für die Sprechabsicht „Vorlieben ausdrücken" (GER A1/A2) unverzichtbar ist.

\textbf{Bezug zur Bedingungsanalyse:}

Die SuS haben das Themenfeld „Sich vorstellen" abgeschlossen. Vorlieben und Hobbys sind ein \textbf{natürlicher Anschluss}: Nach „Wer bin ich?" folgt „Was mag ich?". Dies entspricht auch den Interessen der Lernenden (Freizeit, Musik, Reisen).

\textbf{Fachdidaktische Begründung:}

Die kontrastive Einführung über das deutsche „gefallen" nutzt die Erstsprache als Brücke (positive Transfer-Nutzung). Die visuelle Darstellung der Satzstruktur (Diagramm) unterstützt das Verständnis der ungewohnten Syntax (vgl. Meißner \& Reinfried, 1998, zur Mehrsprachigkeitsdidaktik).

\subsection{Didaktische Reduktion}

\textbf{Was wird NICHT thematisiert:}
\begin{itemize}[leftmargin=2em]
    \item Die Position von \es{a mí, a ti...} vor vs. nach dem Verb (beides akzeptiert)
    \item Nuancen zwischen \es{gustar} und \es{caer bien} (für Personen)
    \item Andere Tempusformen von \es{gustar} (nur Präsens)
    \item Verben mit ähnlicher Struktur (erst Stunde 2)
\end{itemize}

\textbf{Fokus der Stunde:}
\begin{itemize}[leftmargin=2em]
    \item Struktur: Pronomen + \es{gusta/gustan} + Subjekt
    \item Regel: \es{gusta} + Singular/Infinitiv, \es{gustan} + Plural
    \item Wortschatz: 10 Hobbys/Freizeitaktivitäten
\end{itemize}

\subsection{Strukturierung des Unterrichts}

Die Stunde folgt einem \textbf{induktiv-kommunikativen Ansatz}:

\begin{enumerate}[leftmargin=2em]
    \item \textbf{Einstieg (5 min):} Aktivierung durch „Überraschungsfakt" (Problemstellung)
    \item \textbf{Erarbeitung I (10 min):} Struktur entdecken (Kontrastierung DE/ES)
    \item \textbf{Übung I (7 min):} Verständnissicherung + Pattern Drill (isoliert)
    \item \textbf{Erarbeitung II (5 min):} Wortschatz Hobbys
    \item \textbf{Übung II (10 min):} Anwendung (Übersetzen, Partnerarbeit)
    \item \textbf{Sicherung (8 min):} Test + freies Schreiben
\end{enumerate}

\textbf{Alternative:} Bei Zeitproblemen kann der Pattern-Drill (Schritt 5b im Lernpfad) gekürzt werden. Bei schneller Lerngruppe kann die Partnerarbeit um ein Interview erweitert werden.

% ═══════════════════════════════════════════════════════════════════════════
% 6. STUNDENZIELE
% ═══════════════════════════════════════════════════════════════════════════

\section{Stundenziele}

\subsection{Grobziele}

\begin{enumerate}[leftmargin=2em]
    \item Die Schülerinnen und Schüler \textbf{verstehen} die Struktur des Verbs \es{gustar} als Entsprechung zum deutschen „gefallen".
    \item Die Schülerinnen und Schüler \textbf{können} einfache Sätze über Vorlieben bilden und dabei \es{gusta/gustan} korrekt verwenden.
\end{enumerate}

\subsection{Feinziele (operationalisiert)}

\begin{center}
\small
\begin{tabularx}{\textwidth}{|c|X|c|c|}
    \hline
    \textbf{Nr.} & \textbf{Die SuS...} & \textbf{Kompetenz} & \textbf{AFB} \\
    \hline
    1 & ...benennen die indirekten Objektpronomen (me, te, le, nos, os, les). & Sprachliche Mittel & I \\
    \hline
    2 & ...ordnen \es{gusta} (Singular/Infinitiv) und \es{gustan} (Plural) korrekt zu. & Sprachliche Mittel & I \\
    \hline
    3 & ...erklären, warum \es{gustar} wie „gefallen" und nicht wie „mögen" funktioniert. & Sprachbewusstheit & II \\
    \hline
    4 & ...übersetzen einfache Sätze vom Deutschen ins Spanische unter Verwendung von \es{gustar}. & FKK Schreiben & II \\
    \hline
    5 & ...begründen, warum typische Fehler (z.\,B. „*Yo gusto...") falsch sind. & Sprachbewusstheit & III \\
    \hline
    6 & ...verfassen drei eigene Sätze über ihre Vorlieben. & FKK Schreiben & III \\
    \hline
\end{tabularx}
\end{center}

% ═══════════════════════════════════════════════════════════════════════════
% 7. METHODISCHE ÜBERLEGUNGEN
% ═══════════════════════════════════════════════════════════════════════════

\section{Methodische Überlegungen}

\subsection{Einstieg: Überraschung als Aktivierung}

Der Einstieg erfolgt durch einen \textbf{„Überraschungsfakt"}:

\begin{infobox}
    \textit{„Wusstest du? Auf Spanisch kannst du NICHT sagen: ‚Ich mag dich.' \\
    Stattdessen sagst du: ‚Me gustas.' -- wörtlich: ‚Du gefällst mir.'"}
\end{infobox}

\textbf{Begründung:} Diese kognitive Dissonanz weckt Neugier und fokussiert die Aufmerksamkeit auf das Problem: Warum funktioniert \es{gustar} anders als erwartet? Die affektive Dimension (Liebe, Zuneigung) macht das Thema zusätzlich relevant.

\textbf{Alternative:} Ein kurzes Video-Snippet (Telenovela-Szene mit „Me gustas") könnte den Einstieg noch authentischer gestalten, erfordert jedoch mehr Zeit.

\subsection{Erarbeitung: Kontrastierung und Visualisierung}

Die Struktur wird durch \textbf{Kontrastierung} mit dem Deutschen erarbeitet:

\begin{center}
\begin{tabular}{|l|l|}
    \hline
    \textbf{Deutsch (gefallen)} & \textbf{Spanisch (gustar)} \\
    \hline
    \underline{Mir} gefällt \underline{die Musik}. & \underline{Me} gusta \underline{la música}. \\
    \hline
\end{tabular}
\end{center}

Ein \textbf{SVG-Diagramm} im Lernpfad visualisiert die Satzglieder und zeigt mit einem Pfeil: Das Subjekt (la música) bestimmt die Verbform.

\textbf{Begründung:} Die visuelle Darstellung unterstützt das Verständnis der ungewohnten Struktur (Prinzip der multiplen Repräsentation). Die Kontrastierung mit „gefallen" nutzt Vorwissen (Scaffolding).

\textbf{Alternative:} Bei einer rein analogen Stunde könnte ein Tafelanschrieb mit farbiger Kreide dieselbe Funktion übernehmen.

\subsection{Übungsphasen: Vom Drill zur Produktion}

Die Übungen sind \textbf{gestuft} (Scaffolding):

\begin{enumerate}[leftmargin=2em]
    \item \textbf{Stufe 1 (AFB I):} Pronomen einsetzen (isolierte Teilkompetenz)
    \item \textbf{Stufe 2 (AFB I):} gusta/gustan wählen (isolierte Teilkompetenz)
    \item \textbf{Stufe 3 (AFB II):} Sätze übersetzen (Integration beider Teilkompetenzen)
    \item \textbf{Stufe 4 (AFB III):} Freie Produktion (eigene Sätze)
\end{enumerate}

\textbf{Begründung:} Durch das schrittweise Erhöhen der Komplexität können alle SuS Erfolgserlebnisse haben. Leistungsstärkere SuS erreichen Stufe 4, während Stufe 1-2 das Basisniveau sichert.

\subsection{Sozialformen}

\begin{center}
\begin{tabular}{|l|l|l|}
    \hline
    \textbf{Phase} & \textbf{Sozialform} & \textbf{Begründung} \\
    \hline
    Einstieg & Plenum & Gemeinsame Problemstellung \\
    \hline
    Erarbeitung & Einzelarbeit (Lernpfad) & Individuelles Tempo, Sofort-Feedback \\
    \hline
    Übung I & Einzelarbeit & Konzentriertes Üben \\
    \hline
    Übung II & Partnerarbeit & Kommunikative Anwendung, Peer-Korrektur \\
    \hline
    Sicherung & Einzelarbeit & Individuelle Lernstandserhebung \\
    \hline
\end{tabular}
\end{center}

\subsection{Medien}

\begin{itemize}[leftmargin=2em]
    \item \textbf{Digitaler Lernpfad} (HTML): Interaktive Übungen mit Sofort-Feedback
    \item \textbf{Arbeitsblatt} (1 Seite, A4): Grammatik-Kasten, Vokabeln, Übungen zum Eintragen (Hefter)
    \item \textbf{Beamer/Whiteboard}: Für gemeinsame Besprechung des Einstiegs und der Ergebnisse
\end{itemize}

\textbf{Begründung:} Die Kombination aus digitalem Lernpfad (individualisiert, interaktiv) und analogem Arbeitsblatt (zum Abheften, Nachlesen) berücksichtigt verschiedene Lernpräferenzen und schafft ein dauerhaftes Lernprodukt.

\subsection{Differenzierung}

\begin{center}
\begin{tabular}{|l|l|l|}
    \hline
    \textbf{Niveau} & \textbf{Aufgaben} & \textbf{Charakteristik} \\
    \hline
    Basis & AB Ü1, Ü2; Lernpfad Schritt 3, 5 & Reproduktion, Zuordnung \\
    \hline
    Standard & AB Ü3; Lernpfad Schritt 7 & Transfer, Übersetzung \\
    \hline
    Erweitert & AB Ü4, Freies Schreiben; Lernpfad Schritt 8 & Begründung, Produktion \\
    \hline
\end{tabular}
\end{center}

\textbf{S1 (LRS):} Erhält AB in Schriftgröße 12pt, mehr Zeit bei schriftlichen Aufgaben.

\textbf{S2 (Vorkenntnisse):} Wird als Expertin in der Partnerarbeit eingesetzt; erhält Zusatzaufgabe (weitere Verben: \es{encantar, interesar}).

\subsection{Erwartungshorizonte}

\textbf{Zu Übung 3 (Übersetzen):}
\begin{itemize}[leftmargin=2em]
    \item „Mir gefällt Musik." → \es{Me gusta la música.} (Artikel beachten!)
    \item „Dir gefallen die Filme." → \es{Te gustan las películas.}
    \item „Uns gefällt das Reisen." → \es{Nos gusta viajar.}
\end{itemize}

\textbf{Zu Übung 4 (Fehler erklären):}
\begin{itemize}[leftmargin=2em]
    \item „*Me gustan bailar" → Infinitiv ist Singular, daher \es{gusta}.
    \item „*Yo gusto el chocolate" → \es{gustar} funktioniert wie „gefallen"; „el chocolate" ist das Subjekt.
\end{itemize}

\subsection{Überprüfung der Teilziele}

\begin{center}
\begin{tabular}{|c|l|l|}
    \hline
    \textbf{Feinziel} & \textbf{Überprüfung durch} & \textbf{Zeitpunkt} \\
    \hline
    1, 2 & Lernpfad-Übungen (automatisches Feedback) & Min. 15-22 \\
    \hline
    3 & MC-Frage im Lernpfad (Schritt 3) & Min. 12 \\
    \hline
    4 & AB Übung 3 + Lernpfad Schritt 7 & Min. 25-35 \\
    \hline
    5 & AB Übung 4 (schriftliche Begründung) & Min. 35-40 \\
    \hline
    6 & AB Freies Schreiben (3 Sätze) & Min. 40-45 \\
    \hline
\end{tabular}
\end{center}

% ═══════════════════════════════════════════════════════════════════════════
% 8. AUSFÜHRLICHER VERLAUFSPLAN
% ═══════════════════════════════════════════════════════════════════════════

\section{Ausführlicher Verlaufsplan}

\small
\begin{longtable}{|p{1.2cm}|p{2cm}|p{2.5cm}|p{4.5cm}|p{4cm}|}
    \hline
    \textbf{Zeit} & \textbf{Phase} & \textbf{Sozialform / Medien} & \textbf{Lehrerhandlungen} & \textbf{Erwartete Schülerhandlungen} \\
    \hline
    \endfirsthead
    \hline
    \textbf{Zeit} & \textbf{Phase} & \textbf{Sozialform / Medien} & \textbf{Lehrerhandlungen} & \textbf{Erwartete Schülerhandlungen} \\
    \hline
    \endhead

    0-5 min (5') & Einstieg & PL, Beamer & L zeigt „Überraschungsfakt" am Beamer: „Du kannst NICHT sagen: Ich mag dich." Fragt: „Warum?" & SuS lesen, stellen Hypothesen auf. Mögliche Antworten: „Andere Grammatik?", „Anderes Verb?" \\
    \hline

    5-10 min (5') & Erarbeitung I & EA, Lernpfad (Tablet) & L gibt Arbeitsauftrag: „Öffnet den Lernpfad. Bearbeitet Schritt 1-2." Geht herum, beantwortet Fragen. & SuS lesen Schritt 1-2, entdecken: \es{gustar} = gefallen. Studieren Kontrastierung DE/ES, SVG-Diagramm. \\
    \hline

    10-17 min (7') & Übung I & EA, Lernpfad & L: „Bearbeitet Schritt 3 (Verständnis) und Schritt 5 (Pronomen)." Beobachtet Fortschritt. & SuS beantworten MC-Fragen, setzen Pronomen ein. Erhalten Sofort-Feedback. \\
    \hline

    17-22 min (5') & Erarbeitung II & EA, Lernpfad + AB & L teilt AB aus: „Übertrage die Grammatik in Kasten 1+2. Lerne die Vokabeln (Schritt 6)." & SuS übertragen Merksätze ins AB, lernen Vokabeln (Hobbys). \\
    \hline

    22-27 min (5') & Übung II & EA, Lernpfad & L: „Bearbeitet Schritt 5b (gusta/gustan-Drill)." Optional: Schnelle SuS → Schritt 7. & SuS entscheiden: gusta oder gustan? Automatisches Feedback. \\
    \hline

    27-35 min (8') & Anwendung & PA, AB & L: „Partnerarbeit: Bearbeitet AB Ü1-3 gemeinsam. Korrigiert euch gegenseitig." Geht herum, unterstützt. & SuS setzen Pronomen ein (Ü1), wählen gusta/gustan (Ü2), übersetzen (Ü3). Peer-Korrektur. \\
    \hline

    35-40 min (5') & Sicherung I & EA, Lernpfad & L: „Macht den Abschlusstest (Schritt 8)." & SuS beantworten 6 MC-Fragen. Erhalten Score. \\
    \hline

    40-45 min (5') & Sicherung II & EA, AB & L: „Bearbeitet AB Ü4 (Fehler erklären) und schreibt 3 eigene Sätze über eure Vorlieben." & SuS begründen Fehler schriftlich, verfassen eigene Sätze. \\
    \hline

    (45 min) & Ausblick & PL & L sammelt ein, gibt Ausblick: „Nächste Stunde: Gespräche über Vorlieben führen." & SuS geben AB ab, notieren ggf. HA. \\
    \hline
\end{longtable}
\normalsize

\textbf{Legende:}
\begin{itemize}[nosep, leftmargin=2em]
    \item PL = Plenum, EA = Einzelarbeit, PA = Partnerarbeit
    \item L = Lehrkraft, SuS = Schülerinnen und Schüler
    \item AB = Arbeitsblatt
\end{itemize}

% ═══════════════════════════════════════════════════════════════════════════
% 9. REFLEXION (optional, nach der Stunde)
% ═══════════════════════════════════════════════════════════════════════════

\section{Reflexion (nach der Durchführung)}

\textit{[Dieser Abschnitt wird nach der Stunde ausgefüllt.]}

\subsection{Realisierung der Lernziele}

\begin{itemize}[leftmargin=2em]
    \item Wurden die Grobziele erreicht?
    \item Welche Feinziele wurden von wie vielen SuS erreicht?
    \item Ergebnisse des Abschlusstests (Durchschnitt, Streuung)?
\end{itemize}

\subsection{Wirksamkeit des Vorgehens}

\begin{itemize}[leftmargin=2em]
    \item War das Scaffolding (Stufe 1-4) angemessen?
    \item Funktionierte die Partnerarbeit?
    \item War der Zeitplan realistisch?
\end{itemize}

\subsection{Konsequenzen für die Weiterführung}

\begin{itemize}[leftmargin=2em]
    \item Welche Inhalte müssen in Stunde 2 wiederholt werden?
    \item Gibt es individuelle Förderbedarfe?
\end{itemize}

\subsection{Gewonnene Einsichten}

\begin{itemize}[leftmargin=2em]
    \item Eigenes pädagogisches Verhalten
    \item Eigenes methodisches Vorgehen
    \item Eigener fachlicher Stand
\end{itemize}

% ═══════════════════════════════════════════════════════════════════════════
% 10. ANLAGEN
% ═══════════════════════════════════════════════════════════════════════════

\section{Anlagen}

\begin{enumerate}[leftmargin=2em]
    \item \textbf{Arbeitsblatt} (1 Seite, A4): \texttt{arbeitsblatt-gustos.pdf}
    \item \textbf{Digitaler Lernpfad}: \texttt{lernpfad.html}
    \item \textbf{Tafelbild / Einstiegsfolie} (siehe unten)
    \item \textbf{Erwartungshorizont AB} (siehe Abschnitt 7.6)
\end{enumerate}

\subsection*{Anlage 3: Einstiegsfolie (Beamer)}

\begin{infobox}
    \begin{center}
        {\large\bfseries ¿Sabías que...?}

        \vspace{0.5cm}

        En español, NO puedes decir: \textit{„Ich mag dich."}

        \vspace{0.3cm}

        En cambio dices: \es{Me gustas.}

        \vspace{0.3cm}

        Literalmente: \textit{„Du gefällst mir."}

        \vspace{0.5cm}

        {\normalsize\bfseries ¿Por qué?}
    \end{center}
\end{infobox}

% ═══════════════════════════════════════════════════════════════════════════
% 11. PUBLIKATIONSVERZEICHNIS
% ═══════════════════════════════════════════════════════════════════════════

\section{Publikationsverzeichnis}

\subsection*{Rahmenpläne und Bildungsstandards}

Ministerium für Bildung, Wissenschaft und Kultur Mecklenburg-Vorpommern (2019). \textit{Rahmenplan Spanisch für die Qualifikationsphase der gymnasialen Oberstufe}. Schwerin.

Kultusministerkonferenz (2012). \textit{Bildungsstandards für die fortgeführte Fremdsprache (Englisch/Französisch) für die Allgemeine Hochschulreife}. Bonn.

Europarat (2001). \textit{Gemeinsamer europäischer Referenzrahmen für Sprachen: lernen, lehren, beurteilen}. Berlin: Langenscheidt.

\subsection*{Fachwissenschaftliche Literatur}

Bosque, I. \& Demonte, V. (Hrsg.) (1999). \textit{Gramática descriptiva de la lengua española} (Bd. 2). Madrid: Espasa Calpe.

Real Academia Española (2010). \textit{Nueva gramática de la lengua española. Manual}. Madrid: Espasa.

\subsection*{Fachdidaktische Literatur}

Meißner, F.-J. \& Reinfried, M. (Hrsg.) (1998). \textit{Mehrsprachigkeitsdidaktik. Konzepte, Analysen, Lehrerfahrungen mit romanischen Fremdsprachen}. Tübingen: Narr.

Bär, M. (2009). \textit{Förderung von Mehrsprachigkeit und Lernkompetenz}. Tübingen: Narr.

Decke-Cornill, H. \& Küster, L. (2015). \textit{Fremdsprachendidaktik} (3. Aufl.). Tübingen: Narr.

\subsection*{Unterrichtsmaterialien}

Digitaler Lernpfad „Me gusta" (Eigenentwicklung, 2024).

Arbeitsblatt „Me gusta -- Hobbys und Vorlieben" (Eigenentwicklung, 2024).

\vfill

\hrule
\vspace{0.3cm}
\begin{center}
    {\small\color{grau} Unterrichtslangentwurf erstellt gemäß der „Fächerübergreifenden Handreichung zur Erstellung von Unterrichtsentwürfen" (ZLB/IQ M-V)}
\end{center}

\end{document}
